\section{An exact variational principle at the critical exponent}
\label{sec:critical-variational}

For completeness we reproduce the finite-domain variational foundation.
Here $\phi=\tanh$. For arbitrary fixed depth, the critical prior admits a direct finite-domain
large-deviation description. This result requires neither a propagation-of-chaos
assumption nor differentiability of the activation. It identifies the functional used in the critical-rate capacity constructions
with the actual Gaussian probability rate.
Covariance-process large deviations for Gaussian networks have been studied
in a more general functional setting by Andreis, Bassetti, and
Hirsch~\cite{andreis2026}. No novelty claim is made for the general
large-deviation framework; the finite-domain proof below fixes the precise
normalizations and strict-margin convention used here.

Fix $\gamma=1$, enumerate $\mathcal X=\{x^1,\ldots,x^m\}$, where $m=2^d$,
and put $r=L-1$. All limits in this section hold with $d$ and $L$ fixed.
Write $E=[-1,1]^m$ and let $\mathcal P(E)$ carry its weak topology.
For $\nu\in\mathcal P(E)$ define
\[
 K(\nu)=\int_E hh^{\mathsf T}\,d\nu(h),\qquad
 Q_0(a,b)=\frac{\bar x^a\cdot\bar x^b}{d}.
\]
For any positive semidefinite $m\times m$ matrix $Q$, let
\[
 \mathsf T_Q(u)=\bigl(\phi((Q^{1/2}u)_a)\bigr)_{a=1}^m,
 \qquad
 \mathsf M_Q=(\mathsf T_Q)_\#\mathcal N(0,I_m).
\]
Here $T_\#\rho$ is the distribution of $T(U)$ when $U$ has distribution $\rho$,
and $Q^{1/2}$ is the positive semidefinite square root. Relative entropy is
\[
 D(\nu\Vert\rho)=
 \begin{cases}
 \displaystyle\int\log\frac{d\nu}{d\rho}\,d\nu,&\nu\ll\rho,\\
 +\infty,&\text{otherwise}.
 \end{cases}
\]
Finally define the Gaussian energy
\[
 \mathcal G_Q(z)=
 \begin{cases}
 \frac12 z^{\mathsf T}Q^+z,&z\in\operatorname{range}(Q),\\
 +\infty,&z\notin\operatorname{range}(Q),
 \end{cases}
\]
where $Q^+$ is the Moore--Penrose inverse.

\begin{theorem}[Critical finite-domain large deviations]
\label{thm:critical-ldp}
Let $Z_N=(z_w(x^a))_{a=1}^m$ under the unconditioned prior
$\mu_{N,1}$. The sequence $Z_N$ satisfies a large-deviation principle
with speed $N$ and good rate function
\begin{equation}\label{eq:critical-rate}
 \begin{split}
 I(z)=\inf_{\nu_1,\ldots,\nu_r\in\mathcal P(E)}
 \bigg\{&D(\nu_1\Vert\mathsf M_{Q_0})
  +\sum_{s=2}^{r}D\bigl(\nu_s\Vert\mathsf M_{K(\nu_{s-1})}\bigr)
  +\mathcal G_{K(\nu_r)}(z)\bigg\}.
 \end{split}
\end{equation}
In particular, for every open $O\subseteq\mathbb R^m$ and closed
$F\subseteq\mathbb R^m$,
\[
 -\inf_{z\in O}I(z)
 \leq\liminf_{N\to\infty}\frac1N\log\Prb(Z_N\in O),
 \qquad
 \limsup_{N\to\infty}\frac1N\log\Prb(Z_N\in F)
 \leq-\inf_{z\in F}I(z).
\]
The infimum in \eqref{eq:critical-rate} is attained whenever it is finite.
No covariance matrix in the statement is required to be nonsingular.
\end{theorem}

\begin{proof}
We give an exact representation of the finite-width law and then apply
Sanov's theorem and continuous contraction~\cite{dembozeitouni}.

Let $\Gamma_m=\mathcal N(0,I_m)$. Independently for each
$s=1,\ldots,r$, let $U_{s,1},\ldots,U_{s,N}$ have law $\Gamma_m$, and set
\[
 \eta_{s,N}=\frac1N\sum_{i=1}^N\delta_{U_{s,i}}.
\]
Also let $V_N=G/\sqrt N$, where $G\sim\Gamma_m$ is independent of these
empirical measures. Recursively define
\[
 \nu_{1,N}=(\mathsf T_{Q_0})_\#\eta_{1,N},\qquad
 \nu_{s,N}=(\mathsf T_{K(\nu_{s-1,N})})_\#\eta_{s,N}\quad(2\leq s\leq r).
\]
These are exactly the joint laws of the network's empirical hidden-activation
vectors on $\mathcal X$. Indeed, first-layer preactivation vectors are
independent Gaussians with covariance $Q_0$. Conditional on the preceding
layer, each subsequent row of weights gives an independent Gaussian vector
with covariance
\[
 \frac1N\sum_{i=1}^N
 \bigl(h_i^{s-1}(x^a)h_i^{s-1}(x^b)\bigr)_{a,b}
 =K(\nu_{s-1,N}).
\]
The conditional distribution depends on the previous layer only through its
empirical measure. Finally, readout weights have variance $N^{-2}$, so
\begin{equation}\label{eq:critical-output-representation}
 Z_N\ \stackrel{\mathrm{law}}{=}\ K(\nu_{r,N})^{1/2}V_N.
\end{equation}
This representation is exact for every $N$.

Sanov's theorem on $\mathcal P(\mathbb R^m)$ gives the good rate
$D(\eta\Vert\Gamma_m)$ for each $\eta_{s,N}$. The elementary Gaussian
large-deviation principle gives rate $\|v\|_2^2/2$ for $V_N$.
Independence and the finite-product principle therefore give the good rate
\[
 \sum_{s=1}^r D(\eta_s\Vert\Gamma_m)+\frac12\|v\|_2^2
\]
for $(\eta_{1,N},\ldots,\eta_{r,N},V_N)$.

The recursive map from these variables to $(\nu_1,\ldots,\nu_r,z)$ is
continuous. To verify the only potentially delicate point, suppose
$Q_j\to Q$ and $\eta_j\Rightarrow\eta$ weakly. For every continuous
$b:E\to\mathbb R$, the functions $b\circ\mathsf T_{Q_j}$ converge
uniformly on each compact subset of $\mathbb R^m$ to
$b\circ\mathsf T_Q$, since matrix square roots are continuous and $\phi$
is continuous. Their absolute values are uniformly bounded. Tightness of a
weakly convergent sequence of probability measures then gives
\[
 \int b\circ\mathsf T_{Q_j}\,d\eta_j
 \longrightarrow\int b\circ\mathsf T_Q\,d\eta.
\]
Also $\nu\mapsto K(\nu)$ is continuous because every coordinate
$h_ah_b$ is continuous and bounded on $E$. This proves continuity of the
whole recursive map, including \eqref{eq:critical-output-representation}.

Continuous contraction now gives the rate by minimizing over its preimages.
For a measurable map $T$ and a reference probability $\rho$, the entropy
contraction identity is
\begin{equation}\label{eq:entropy-contraction-identity}
 \inf_{\eta:T_\#\eta=\nu}D(\eta\Vert\rho)
 =D(\nu\Vert T_\#\rho).
\end{equation}
The lower bound is data processing. When the right side is finite and
$q=d\nu/d(T_\#\rho)$, the measure $d\eta=q(T(u))\,d\rho(u)$ attains
equality. If the right side is infinite, data processing already proves
the identity in the extended real sense. Apply
\eqref{eq:entropy-contraction-identity} successively with
$T=\mathsf T_{Q_0}$ and $T=\mathsf T_{K(\nu_{s-1})}$.
The last minimization is
\[
 \inf_{v:Q^{1/2}v=z}\frac12\|v\|_2^2=\mathcal G_Q(z),
\]
including singular $Q$. This proves \eqref{eq:critical-rate}.

The contraction theorem supplies a good rate and exponential tightness.
For completeness, the displayed minimization is also directly well behaved:
relative entropy is jointly lower semicontinuous, the map
$\nu\mapsto\mathsf M_{K(\nu)}$ is weakly continuous, and
\[
 \mathcal G_Q(z)=\sup_{t\in\mathbb R^m}
 \left\{t^{\mathsf T}z-\frac12t^{\mathsf T}Qt\right\}
\]
is jointly lower semicontinuous in $(Q,z)$. Since $\mathcal P(E)^r$ is
compact, every finite infimum in \eqref{eq:critical-rate} is attained.
Moreover $K(\nu)\preceq mI_m$, so $I(z)\geq\|z\|_2^2/(2m)$.
\end{proof}

\begin{corollary}[Full-margin masses and posterior rates]
\label{cor:critical-margin-rates}
For a Boolean function $f$, regarded as a vector in $\{-1,1\}^m$, set
\[
 J(f)=\inf_{z:\ f_a z_a\geq1/2\ \forall a}I(z).
\]
Then the strict-margin event used throughout this paper has the exact rate
\begin{equation}\label{eq:critical-full-margin-rate}
 \lim_{N\to\infty}-\frac1N\log\mu_{N,1}(A_{N,f})=J(f).
\end{equation}
For a fixed noiseless training set $S$, let $\mathcal B_S$ be the Boolean
functions agreeing with its labels. If
$J_S=\min_{g\in\mathcal B_S}J(g)<\infty$, the actual globally conditioned
posterior satisfies, for every $f\in\mathcal B_S$,
\[
 \lim_{N\to\infty}-\frac1N\log\Pi_{N,1}(f\mid S)=J(f)-J_S.
\]
In particular, posterior mass concentrates on the minimizers of $J$ within
$\mathcal B_S$. These logarithmic asymptotics do not determine how mass is
divided between tied minimizers.
\end{corollary}

\begin{proof}
For $c\geq1$, retaining any feature laws in \eqref{eq:critical-rate} and
scaling only the output gives
\[
 I(cz)\leq c^2 I(z),
\]
because the entropy terms are nonnegative and the Gaussian energy scales
quadratically. If $f_a z_a\geq1/2$ for all $a$, then $cz$ has strict
margin for every $c>1$. Letting $c\downarrow1$ shows that the infimum of
$I$ over the open strict-margin orthant equals its infimum over the
closed orthant. The lower and upper large-deviation bounds therefore agree,
proving \eqref{eq:critical-full-margin-rate}.

Conditioning the prior on the global margin event and subsequently on the
labels gives exactly
\[
 \Pi_{N,1}(f\mid S)=
 \frac{\mu_{N,1}(A_{N,f})}
 {\sum_{g\in\mathcal B_S}\mu_{N,1}(A_{N,g})}.
\]
For fixed $d$ this is a finite sum, whose exponential rate is the smallest
$J(g)$ among its terms. This proves the ratio formula and concentration
claim.
\end{proof}

